\documentclass[12pt]{article}
\usepackage[utf8]{inputenc}
\usepackage[T1]{fontenc}
\usepackage[margin=1in]{geometry}
\usepackage{setspace}
\usepackage{xcolor}
\usepackage[colorlinks=true, urlcolor=blue!60!black, linkcolor=blue!60!black]{hyperref}
\setlength{\parindent}{0em}
\setlength{\parskip}{0.6em}
\renewcommand{\baselinestretch}{1.0}
\begin{document}
\begin{center}
{\LARGE \textit{Virginia v. Black}, 538 U.S. 343 (2003)\par}
\vspace{0.5em}
{\large Case Brief --- Two-Page Edition\par}
\end{center}

\section{Memory Jogger}
A State may ban cross burning with intent to intimidate, a form of true threat, without violating \textit{R.A.V.}; a four-Justice plurality found Virginia's prima facie evidence provision invalid as construed in Barry Black's jury instruction (\href{https://supreme.justia.com/cases/federal/us/538/343/}{538 U.S. 343}).

\section{Detailed Case Facts}
Virginia made it a felony to burn a cross with intent to intimidate and treated any cross burning as ``prima facie evidence of an intent to intimidate.'' Barry Black led a Klan rally on private property that ended in a cross burning; his jury was told the burning alone sufficed to infer intent. Richard Elliott and Jonathan O'Mara tried to burn a cross on an African-American neighbor's lawn.

\section{Procedural History}
The Supreme Court of Virginia struck the statute down under \textit{R.A.V.} as content discrimination and as overbroad. The U.S. Supreme Court granted certiorari.

\section{Judicial Votes}
Affirmed as to Black (7--2), vacated and remanded as to Elliott and O'Mara (5--4). O'Connor wrote for the Court in Parts I--III and for a plurality of four in Parts IV--V; Stevens concurred; Scalia and Souter (with Kennedy and Ginsburg) concurred in part and dissented in part; Thomas dissented.

\section{Holding}
A ban on cross burning with intent to intimidate is valid because such intimidation is a true threat and cross burning a ``particularly virulent form of intimidation.'' The plurality held that a prima facie provision permitting conviction from the act alone is facially invalid.

\section{Analysis of Opinions}
O'Connor traced cross burning's ties to Klan terror, defined intimidation as a true threat meant to cause fear of bodily harm, and held that \textit{R.A.V.} allows singling out the most virulent threats when the ban applies regardless of motive or victim. The plurality found that the prima facie clause ignores context, chilling political and artistic uses. Scalia thought it not substantially overbroad; Souter saw viewpoint suppression; Thomas called cross burning ``conduct, not expression.'' \textit{Counterman v. Colorado} later required proof of at least recklessness.

\section{Examples: Future Applications}
\textbf{Same-Side:} a ban on hanging a noose to intimidate supports a conviction on independent proof of intent. \textbf{Opposite-Side:} convicting a protest artist on a prima facie clause alone is invalid. \textbf{Fence-Sitter:} posting a burning-cross image that tags a Black candidate.

\section{Critique}
Souter's critique asks whether a ban on a symbol fused with one ideology is really viewpoint neutral, and Scalia's asks whether a jury instruction can make a valid statute facially void.

\section{Key Quotations}
\begin{quote}
``\,`True threats' encompass those statements where the speaker means to communicate a serious expression of an intent to commit an act of unlawful violence to a particular individual or group of individuals. \dots Intimidation in the constitutionally proscribable sense of the word is a type of true threat, where a speaker directs a threat to a person or group of persons with the intent of placing the victim in fear of bodily harm or death.''
\end{quote}
\noindent --- Justice O'Connor, Opinion of the Court, 538 U.S. at 359--360.

\end{document}
